% Independent fifth-power contribution. Standard amsmath/amsthm only.
% Labels have prefix q5:. All finite parts use the coordinate x itself.
\section{The fifth digamma power and its mixed coordinate}
\label{q5:section}

This section resolves the fifth-power calculation posed in the inspected
quartic report.  It fixes all endpoint constants and identifies the
additional mixed spectral combination which occurs beside the
height-one Laurent coefficients.  The result does not assert that this
mixed combination is arithmetically independent of the height-one data.

Write $z_j=\zeta(j)$ and $\gamma=\gamma_0$.  Strict multiple zeta
functions have decreasing indices:
\[
 Z_r(s_1,\ldots,s_r)
   =\sum_{n_1>\cdots>n_r\geq1}
        n_1^{-s_1}\cdots n_r^{-s_r}.
\]
For their height-one germs put
\begin{equation}\label{q5:lambda}
 \lambda_r=[u^1]Z_r(1+u,\underbrace{1,\ldots,1}_{r-1}),
 \qquad
 \lambda_1=-\gamma_1,\quad
 \lambda_2=\eta,\quad\lambda_3=\theta,\quad\lambda_4=\xi.
\end{equation}
The bracket extracts the coefficient itself, with no factorial or sign.
The coordinate finite part of an integral is the constant coefficient
of its cutoff expansion in $\varepsilon$ for the integral from
$\varepsilon$ to the stated upper endpoint.

\subsection{The identities}

Define
\begin{align}
 R(s)&=Z_2(s,3)-z_3\zeta(s),\label{q5:R}\\
 J(s)&=2R(s)+2Z_2(s+2,1)+Z_2(s+1,2)
                 -2\gamma\zeta(s+2)+z_2\zeta(s+1).
 \label{q5:J}
\end{align}
Here $R$ is holomorphic in $\operatorname{Re}s>-1$, and $J$ is
holomorphic in $\operatorname{Re}s>0$.  In particular $J'(1)$ is an
ordinary derivative of a holomorphic function, not a coefficient
requiring a new choice of finite-part normalization.

\begin{theorem}[Fifth-power and mixed-polygamma bridges]
\label{q5:main}
Set
\[
 Q_5=\operatorname{FP}\int_0^1\psi(x)^5\,dx,
 \qquad
 I=\operatorname{FP}\int_0^1\psi(x)\psi'(x)^2\,dx.
\]
Then
\begin{equation}\label{q5:Ibridge}
 \boxed{\displaystyle
 I=-2\gamma z_3+\frac{35}{6}z_4+2J'(1).}
\end{equation}
The fifth-power identity is
\begin{equation}\label{q5:bridge}
 \boxed{\begin{aligned}
 Q_5={}&30\gamma^2z_2-60\gamma z_3+\frac{35}{3}z_4
       -120\xi+120\gamma\theta\\
 &-60(\gamma^2-z_2)\eta
       +20(\gamma^3-3\gamma z_2+2z_3)\lambda_1
       +10J'(1).
 \end{aligned}}
\end{equation}
Equivalently, the mixed integral itself isolates the extra coordinate:
\begin{equation}\label{q5:QminusI}
 \boxed{\begin{aligned}
 Q_5-5I={}&30\gamma^2z_2-50\gamma z_3-\frac{35}{2}z_4
       -120\xi+120\gamma\theta\\
 &-60(\gamma^2-z_2)\eta
       +20(\gamma^3-3\gamma z_2+2z_3)\lambda_1.
 \end{aligned}}
\end{equation}
All the ordinary zeta derivatives in the compact notation are explicit:
\begin{align}\label{q5:Jexpanded}
 J'(1)={}&2R'(1)
       +2\left.\frac{\partial}{\partial s}Z_2(s,1)\right|_{s=3}
       +\left.\frac{\partial}{\partial s}Z_2(s,2)\right|_{s=2}
       -2\gamma\zeta'(3)+z_2\zeta'(2).
\end{align}
If $\phi_{13}=[u^1]Z_2(1+u,3)$, then
$R'(1)=\phi_{13}+z_3\gamma_1=\phi_{13}-z_3\lambda_1$.
\end{theorem}

The last conversion matters.  Substituting $\phi_{13}$ for $R'(1)$
without its $z_3\gamma_1$ correction would give an incorrect fifth-power
formula.  The mixed term is one specified linear combination.  No claim
that its separate components form a minimal arithmetic basis is needed.

\subsection{A convergent arithmetic form and the spectral germs}

For $n\geq1$ use
\begin{equation}\label{q5:integer-bcd}
 b=\psi(n),\qquad c=2z_2-\psi'(n),\qquad
 d=\frac12\psi''(n),\qquad L_n=\log n,
\end{equation}
where $b,c,d$ in a summand always refer to that summand's $n$.
Then
\begin{align}
 J(s)&=\sum_{n\geq1}\frac1{n^s}
      \left(2d+\frac{c}{n}+\frac{2b}{n^2}\right),
      &&\operatorname{Re}s>0,\label{q5:Jseries}\\
 I&=-2\gamma z_3+\frac{35}{6}z_4
       -\sum_{n\geq1}
        \left(\frac{4d}{n}+\frac{2c}{n^2}
                         +\frac{4b}{n^3}\right)L_n.
 \label{q5:Iarithmetic}
\end{align}
Another expression for $Q_5$ is
\begin{align}\label{q5:Qarithmetic}
 Q_5={}&30\gamma^2z_2-60\gamma z_3+\frac{35}{3}z_4
           +20\gamma_4-120z_2\gamma_2\notag\\
 &+20\sum_{n\geq1}
     \frac{(b^3-3bc+d-L_n^3+6z_2L_n)L_n}{n}\notag\\
 &-20\sum_{n\geq1}\frac{bL_n}{n^3}
       -10\sum_{n\geq1}\frac{cL_n}{n^2}.
\end{align}
Every series in these displays converges absolutely.  Indeed
$\psi(n)=\log n+O(n^{-1})$,
$\psi'(n)=O(n^{-1})$, and $\psi''(n)=O(n^{-2})$.
The summand in the first sum of \eqref{q5:Qarithmetic} is
$O((1+\log^3n)/n^2)$.  These estimates also prove normal convergence
of \eqref{q5:Jseries} and of each fixed spectral derivative on compact
subsets of its half-plane.  The series for $R$ is
\[
 R(s)=\sum_{n\geq1}
       \frac{H_{n-1}^{(3)}-z_3}{n^s};
\]
its summands are $O(n^{-\operatorname{Re}s-2})$.

The height-one Laurent germs used here have a direct existence proof.
For $r\geq1$ and $\operatorname{Re}s>1$,
\begin{equation}\label{q5:mellin}
 Z_r(s,1^{r-1})=
 \frac{(-1)^{r-1}}{(r-1)!\,\Gamma(s)}
 \int_0^\infty
    \frac{t^{s-1}\log^{r-1}(1-e^{-t})}{e^t-1}\,dt.
\end{equation}
This follows from the elementary-symmetric generating series and
termwise Mellin integration.  Near zero the kernel is its displayed
leading term
$(-1)^{r-1}\log^{r-1}t/((r-1)!t)$ plus
$O(1+|\log t|^{r-1})$.  Subtraction of that term gives
\[
 \Gamma(1+u)Z_r(1+u,1^{r-1})=u^{-r}+V_r(u),
 \qquad V_r(0)=0,
\]
with $V_r$ holomorphic near zero.  The assertion $V_r(0)=0$ follows
by integrating the derivative of $\log^r(1-e^{-t})$ and comparing its
lower boundary with $\log^r t$.  Thus the principal and constant
coefficients are fixed by $1/\Gamma(1+u)$, and in particular $\xi$
in \eqref{q5:lambda} is well defined.

\subsection{Global partial fractions with a fixed entire remainder}

For $n\geq0$, set
\begin{equation}\label{q5:localcoeffs}
 \begin{aligned}
 b_n&=H_n-\gamma,&
 c_n&=z_2+H_n^{(2)},\\
 d_n&=H_n^{(3)}-z_3,&
 e_n&=z_4+H_n^{(4)}.
 \end{aligned}
\end{equation}
The local digamma expansion is
\[
 \psi(-n+\varepsilon)
 =-\varepsilon^{-1}+b_n+c_n\varepsilon
       +d_n\varepsilon^2+e_n\varepsilon^3+O(\varepsilon^4).
\]
For $F_5(z)=\psi(z)^5$ and $F_I(z)=\psi(z)\psi'(z)^2$, let
\[
 P_{X,n}(z)=\sum_{j=1}^5\frac{p_{X,j,n}}{(z+n)^j},
 \qquad X\in\{5,I\}.
\]
The principal coefficients, with subscripts $n$ suppressed in the table,
are
\begin{center}
\begin{tabular}{c|c|c}
$j$&$p_{5,j}$&$p_{I,j}$\\\hline
5&$-1$&$-1$\\
4&$5b$&$b$\\
3&$-10b^2+5c$&$-c$\\
2&$10b^3-20bc+5d$&$2bc-3d$\\
1&$-5b^4+30b^2c-20bd-10c^2+5e$&$4bd+c^2-5e$
\end{tabular}
\end{center}
Write $r_{X,n}=p_{X,1,n}$.  The constant coefficients at zero are
\begin{align}
 \kappa_5={}&-\gamma^5+20\gamma^3z_2-30\gamma^2z_3
         -30\gamma z_2^2+20\gamma z_4+20z_2z_3-5z_5,
 \label{q5:kappa5}\\
 \kappa_I={}&-\gamma z_2^2-6\gamma z_4-2z_2z_3+7z_5.
 \label{q5:kappaI}
\end{align}

\begin{lemma}[Two normalized Mittag--Leffler expansions]
\label{q5:ML}
For $X=5,I$ and $z\notin\{0,-1,-2,\ldots\}$,
\begin{equation}\label{q5:MLformula}
 F_X(z)=P_{X,0}(z)+\kappa_X+
       \sum_{n\geq1}\bigl(P_{X,n}(z)-P_{X,n}(0)\bigr).
\end{equation}
The series converges normally on compact sets avoiding the poles.
\end{lemma}
\begin{proof}
The coefficient table follows by finite multiplication of the local
digamma series and its derivative.  For the entire remainder, put
$f_X=F_X-P_{X,0}$ and integrate
$z f_X(w)/(w(w-z))$ on $|w|=N+1/2$.
Digamma asymptotics in the right half-plane and the reflection formula
in the left half-plane give $f_X(w)=O(1+\log^5N)$ on these circles.
The cotangent and its fixed derivatives are uniformly bounded there:
away from the real axis use their exponential expressions, and on the
remaining arcs the real coordinate stays a fixed distance from the
integers.  The contour integral is therefore
$O((1+\log^5N)/N)$ and tends to zero.
Its residues are $f_X(z)$, $-f_X(0)$, and
$-(P_{X,n}(z)-P_{X,n}(0))$.  This proves the identity.
Finally $p_{5,j,n}=O(1+\log^{5-j}n)$ and the corresponding mixed
coefficients have no greater growth.  The two-point subtractions give
an additional inverse power of $n$, proving normal convergence.
\end{proof}

\subsection{The finite summation ledger}

Integrating \eqref{q5:MLformula} and telescoping each pole of order
at least two gives
\begin{equation}\label{q5:integrated}
 \operatorname{FP}\int_0^1F_X(x)\,dx
   =B_X+\sum_{n\geq1}r_{X,n}
        \left(\log\left(1+\frac1n\right)-\frac1n\right),
\end{equation}
where
\begin{equation}\label{q5:Bmechanism}
 B_X=\kappa_X+
  \sum_{j=2}^5\sum_{n\geq1}
    \left(\frac{p_{X,j,n}-p_{X,j,n-1}}{(j-1)n^{j-1}}
                       -\frac{p_{X,j,n}}{n^j}\right).
\end{equation}
This formula includes the finite parts at the origin.  For example,
the lower endpoint contribution of $p_{X,j,0}x^{-j}$ is
$-p_{X,j,0}/(j-1)$; it cancels the initial term of that pole's finite
telescoping sum.  The simple-pole finite part is zero.

Here is a shorter directly reproducible form of all the nonlogarithmic
terms.  In the sums in this display only, let
$b=b_{n-1}$, $c=c_{n-1}$, $d=d_{n-1}$.  Then
\begin{align}
 B_5={}&\kappa_5+\sum_{n\geq1}\left\{
 \frac{-10b^3+20bc-5d+30b^2-20c}{n^2}
 +\frac{-20b^2+15c}{n^3}
 +\frac{5b-35/6}{n^4}+\frac6{n^5}\right\},
 \label{q5:B5sum}\\
 B_I={}&\kappa_I+\sum_{n\geq1}\left\{
 \frac{-2bc+3d+2c}{n^2}+\frac{2b-c}{n^3}
 +\frac{-3b-7/6}{n^4}+\frac2{n^5}\right\}.
 \label{q5:BIsum}
\end{align}
All these sums converge absolutely.

To evaluate them and the finite residue sums, write
\[
 H_N(a_1,\ldots,a_d)=
    \sum_{N\geq n_1>\cdots>n_d\geq1}
                 n_1^{-a_1}\cdots n_d^{-a_d},
 \qquad e_j(N)=H_N(1^j).
\]
The finite product rule is the quasi-shuffle rule with a positive merged
index.  The summation rule is equally elementary:
\begin{equation}\label{q5:finite-prefix}
 \sum_{n=1}^N\frac{H_n(a_1,\ldots,a_d)}{n^p}
   =H_N(p,a_1,\ldots,a_d)
      +H_N(p+a_1,a_2,\ldots,a_d),
\end{equation}
with the second term absent for an empty word.  Both statements follow
by separating strict inequalities from equalities of the outer indices.
There is no infinite rearrangement in either identity.

For completeness, the convergent values required in this section are
\begin{equation}\label{q5:MZVtable}
\begin{array}{c|c}
\text{value}&\text{reduction}\\\hline
Z_2(2,1)&z_3\\
Z_2(3,1),\ Z_2(2,2),\ Z_3(2,1,1)&z_4/4,\ 3z_4/4,\ z_4\\
Z_2(4,1),\ Z_3(3,1,1)&2z_5-z_2z_3\\
Z_2(3,2),\ Z_3(2,2,1)&3z_2z_3-11z_5/2\\
Z_2(2,3),\ Z_3(2,1,2)&9z_5/2-2z_2z_3\\
Z_4(2,1,1,1)&z_5
\end{array}
\end{equation}
These are the classical low-weight reductions.  Their use here can be
checked from the finite product rule, the depth-two sum theorem and
iterated-integral duality.  At weight five, for example, the stuffle
and shuffle products are respectively
\[
 z_2z_3=Z_2(2,3)+Z_2(3,2)+z_5,
 \qquad
 z_2z_3=Z_2(2,3)+3Z_2(3,2)+6Z_2(4,1).
\]
Together with
$Z_2(4,1)+Z_2(3,2)+Z_2(2,3)=z_5$, these give the three
depth-two weight-five entries.  Substitution $t\mapsto1-t$ in the
iterated integrals gives the paired dual values in the table.
The same arguments at weights three and four give the first two rows;
$z_2^2=5z_4/2$ is used throughout.

Here is an explicit finite identity which records the full fifth-power
constant.  In this display set $h_j=H_N^{(j)}$, $h=h_1$,
$E_j=e_j(N)$, and abbreviate $H_N(a_1,\ldots,a_d)$ by
$H(a_1,\ldots,a_d)$.  Put
\[
 A=\gamma^3-3\gamma z_2+z_3,\qquad
 r_0=-5\gamma^4+30\gamma^2z_2-20\gamma z_3-10z_2^2+5z_4.
\]
Then the exact equality for every $N\geq1$ is
\begin{align}\label{q5:S5finite}
 \sum_{n=1}^N\frac{r_{5,n}}n={}&
  -120E_5+120\gamma E_4+60(z_2-\gamma^2)E_3
       +20AE_2+r_0h+20Ah_2\notag\\
 &-120H(2,1,1,1)+120\gamma H(2,1,1)
       +(50z_2-60\gamma^2)H(2,1)\notag\\
 &+20E_2h_3-20H(3,1,1)-20hh_4
       +20H(4,1)+20h_5\notag\\
 &+10hH(2,2)-10H(2,1,2)-10H(2,2,1)+10H(2,3)\notag\\
 &-20\gamma hh_3+20\gamma H(3,1)+10z_2hh_2.
\end{align}
One derivation first writes the residue polynomial as
\begin{align*}
 r_{5,n}={}&-120e_4(n)+120\gamma e_3(n)
       +60(z_2-\gamma^2)e_2(n)+20A H_n+r_0\\
 &+20(H_n-\gamma)H_n^{(3)}
       +5\bigl(H_n^{(2)}\bigr)^2-25H_n^{(4)}+10z_2H_n^{(2)}.
\end{align*}
Apply \eqref{q5:finite-prefix} to every term, and then use the finite
product rule to collect the products occurring in \eqref{q5:S5finite}.

For the mixed residue there is a shorter identity.  Put
$C_0=4\gamma z_3+z_2^2-5z_4$.  Then
\begin{align}\label{q5:SIfinite}
 \sum_{n=1}^N\frac{r_{I,n}}n={}&
 4E_2(h_3-z_3)
 +h\{2H(2,2)-4h_4-4\gamma h_3+2z_2h_2+C_0\}\notag\\
 &-4H(3,1,1)-2H(2,1,2)-2H(2,2,1)+2H(2,3)\notag\\
 &+4H(4,1)+4h_5+4\gamma H(3,1)-4z_3h_2-2z_2H(2,1).
\end{align}
This is again a finite product-and-prefix calculation.

\begin{lemma}[All finite constants]
\label{q5:constants}
As $N\to\infty$, with $L=\log N$,
\begin{align}
 \sum_{n=1}^N\frac{r_{5,n}}n
   &=-L^5+20z_2L^3-90z_4L+C_5+o(1),
   \label{q5:S5asymptotic}\\
 r_{5,N}\log(N+1)
   &=-5L^5+60z_2L^3-90z_4L+o(1),
   \label{q5:r5asymptotic}
\end{align}
where
\begin{equation}\label{q5:C5}
 C_5=-\gamma^5+30\gamma^3z_2-80\gamma^2z_3
         -\frac{65}{2}\gamma z_4+60z_2z_3-69z_5.
\end{equation}
For the mixed residue,
\begin{equation}\label{q5:CI}
 \sum_{n\geq1}\frac{r_{I,n}}n
      =C_I=\gamma z_4-12z_2z_3+15z_5,
 \qquad r_{I,N}\log(N+1)=o(1).
\end{equation}
Finally the nonlogarithmic terms satisfy
\begin{equation}\label{q5:Bconstants}
 B_5-C_5=30\gamma^2z_2-60\gamma z_3+\frac{35}{3}z_4,
 \qquad
 B_I-C_I=-2\gamma z_3+\frac{35}{6}z_4.
\end{equation}
\end{lemma}
\begin{proof}
In \eqref{q5:S5finite}, substitute
$h=\log N+\gamma+O(N^{-1})$ and
$h_j=z_j+O(N^{1-j})$ for $j\geq2$, and replace the words with
outer exponent at least two by their limits in
\eqref{q5:MZVtable}.  Every discarded term is $o(1)$ even after
multiplication by the displayed powers of $\log N$.
The elementary polynomials are fixed, without any interpolation, by
\[
 E_0=1,\qquad
 jE_j=\sum_{r=1}^j(-1)^{r-1}h_rE_{j-r}.
\]
This gives \eqref{q5:S5asymptotic} including its constant
\eqref{q5:C5}.  For \eqref{q5:r5asymptotic}, directly substitute
$b_N=L+O(N^{-1})$, $c_N=2z_2+O(N^{-1})$,
$d_N=O(N^{-2})$, and $e_N=2z_4+O(N^{-3})$ in the residue.
Its constant polynomial term is
$-40z_2^2+10z_4=-90z_4$.

For \eqref{q5:SIfinite}, the first term tends to zero.  The expression
in braces multiplying $h$ has zero limit by
$Z_2(2,2)=3z_4/4$ and $z_2^2=5z_4/2$; its convergence is fast
enough that its product with $h$ tends to zero.  The remaining convergent
words have value
\begin{align*}
 &-4Z_3(3,1,1)-2Z_3(2,1,2)-2Z_3(2,2,1)+2Z_2(2,3)\\
 &\qquad+4Z_2(4,1)+4z_5+4\gamma Z_2(3,1)-6z_2z_3\\
 &\hspace{15mm}=\gamma z_4-12z_2z_3+15z_5.
\end{align*}
Also $r_{I,n}=O(n^{-1})+O(\log n/n^2)$, since its prospective
constant term $4z_2^2-10z_4$ vanishes.  This proves the ordinary
convergence and the boundary assertion in \eqref{q5:CI}.
Finally, expand the products in \eqref{q5:B5sum} and
\eqref{q5:BIsum} by the finite quasi-shuffle rule, prefix their
outer denominator, and use \eqref{q5:MZVtable}.
The results are respectively $C_5$ and $C_I$ plus the two explicitly
displayed corrections in \eqref{q5:Bconstants}.
\end{proof}

The preceding calculation is a full constant-term computation.  In
particular it retains the $-90z_4\log N$ terms on both sides of
the summation-by-parts boundary; neither term may be dropped before
their cancellation.

\subsection{Identification of the logarithmic terms}

The residue differences are particularly short.  For $n\geq1$,
with $b=b_{n-1}$, $c=c_{n-1}$, $d=d_{n-1}$,
\begin{align}
 r_{5,n}-r_{5,n-1}
    &=-\frac{20(b^3-3bc+d)}n+\frac{20b}{n^3}
                              +\frac{10c}{n^2},
      \label{q5:diff5}\\
 r_{I,n}-r_{I,n-1}
    &=\frac{4d}n+\frac{4b}{n^3}+\frac{2c}{n^2}.
      \label{q5:diffI}
\end{align}
For any residue sequence, finite summation by parts gives
\begin{equation}\label{q5:logSBP}
 \sum_{n=1}^Nr_{X,n}\log\left(1+\frac1n\right)
  =r_{X,N}\log(N+1)
       -\sum_{n=1}^N(r_{X,n}-r_{X,n-1})\log n.
\end{equation}
For $I$, equations \eqref{q5:integrated}, \eqref{q5:CI},
\eqref{q5:Bconstants}, and \eqref{q5:diffI} immediately give
\eqref{q5:Iarithmetic} and hence \eqref{q5:Ibridge}.

For $Q_5$, define
\[
 M_N=\sum_{n=1}^N\frac{(b_{n-1}^3-3b_{n-1}c_{n-1}+d_{n-1})\log n}{n}.
\]
Equations \eqref{q5:S5asymptotic}--\eqref{q5:logSBP} give
\begin{align}\label{q5:limit}
 Q_5={}&B_5-C_5+
       20\lim_{N\to\infty}
          \left(M_N-\frac{\log^5N}{5}+2z_2\log^3N\right)\notag\\
 &-20\sum_{n\geq1}\frac{b_{n-1}\log n}{n^3}
       -10\sum_{n\geq1}\frac{c_{n-1}\log n}{n^2}.
\end{align}
Splitting off $\log^3n-6z_2\log n$ and using the ordinary
Stieltjes limits gives \eqref{q5:Qarithmetic}.

For the exact spectral form introduce, initially for
$\operatorname{Re}s>1$,
\[
 T(s)=\sum_{n\geq1}
       \frac{b_{n-1}^3-3b_{n-1}c_{n-1}+d_{n-1}}{n^s}.
\]
Finite elementary-symmetric algebra gives
\begin{align}\label{q5:Tidentity}
 T(s)={}&6Z_4(s,1,1,1)-6\gamma Z_3(s,1,1)
           +3(\gamma^2-z_2)Z_2(s,1)\notag\\
 &-(\gamma^3-3\gamma z_2+2z_3)\zeta(s)-R(s).
\end{align}
Moreover $T(s)+\zeta'''(s)-6z_2\zeta'(s)$ is holomorphic in
$\operatorname{Re}s>0$.  Thus $T$ has principal part
$6u^{-4}-6z_2u^{-2}$ at $s=1+u$, and the finite limit in
\eqref{q5:limit} is exactly $-[u^1]T(1+u)$.
One can see the sign directly by differentiating the normally convergent
subtracted Dirichlet series; the ordinary limits contribute
$\gamma_4-6z_2\gamma_2$.
By \eqref{q5:Tidentity},
\[
 [u^1]T(1+u)=6\xi-6\gamma\theta+3(\gamma^2-z_2)\eta
       -(\gamma^3-3\gamma z_2+2z_3)\lambda_1-R'(1).
\]
The two remaining sums in \eqref{q5:limit} are
\[
 -\left.\frac{\partial}{\partial s}Z_2(s,1)\right|_{s=3}
          +\gamma\zeta'(3),\qquad
 -\left.\frac{\partial}{\partial s}Z_2(s,2)\right|_{s=2}
          -z_2\zeta'(2),
\]
respectively.  Substitution proves \eqref{q5:bridge}; subtraction
of five times \eqref{q5:Ibridge} gives \eqref{q5:QminusI}.

\subsection{A differential-polynomial interpretation and exact endpoints}

The fifth inverse-Gamma differential polynomial is
\begin{align}\label{q5:A5}
 \mathcal A_5
 ={}&\psi^5-10\psi^3\psi'+15\psi(\psi')^2+10\psi^2\psi''
            -10\psi'\psi''-5\psi\psi'''+\psi''''.
\end{align}
It equals $-\Gamma\,D^5(1/\Gamma)$.
Let $\Delta G=G(1)-[x^0]G(x)$ for a Laurent-meromorphic function
at zero.  Coordinate finite parts obey
$\operatorname{FP}\int_0^1G'(x)\,dx=\Delta G$.
The complete endpoint table needed for \eqref{q5:A5} is
\begin{equation}\label{q5:endpoints}
\begin{array}{c|c}
G&\Delta G\\\hline
\psi^4&12\gamma^2z_2-12\gamma z_3-11z_4\\
\psi^2\psi'&2\gamma z_3+3z_4/2\\
(\psi')^2&-6z_4\\
\psi\psi''&8z_4\\
\psi'''&0
\end{array}
\end{equation}
Integration by parts, with these constants retained, gives
\begin{align}\label{q5:Bell-explanation}
 Q_5-5I
   &=\operatorname{FP}\int_0^1\mathcal A_5(x)\,dx
       +30\gamma^2z_2-50\gamma z_3-\frac{35}{2}z_4.
\end{align}
Consequently \eqref{q5:QminusI} is compatible with the all-degree
inverse-Gamma construction, while the raw fifth power contains $I$.
This distinguishes a proved differential-polynomial cancellation from
the separate proposed cancellation by a scalar polynomial in $\psi$.
The formulas here neither prove nor disprove the existence of a further
ordinary-constant identity eliminating $J'(1)$.

\subsection{Normalized primitives and the polygamma hierarchy}

For $n\geq1$, define on $\operatorname{Re}z>0$
\[
 F_1(n,z)=\log\left(1+\frac zn\right)-\frac zn,
 \qquad
 F_j(n,z)=\frac{(n+z)^{1-j}-n^{1-j}}{1-j}-\frac z{n^j}
 \quad(j\geq2),
\]
with the logarithm real for positive arguments.

\begin{corollary}[Two normalized primitives]\label{q5:primitives}
For $X\in\{5,I\}$ the expression
\begin{align}\label{q5:primitiveformula}
 \mathcal P_X(z)={}&r_{X,0}\log z
       +\sum_{j=2}^5\frac{p_{X,j,0}z^{1-j}}{1-j}
       +\kappa_Xz
       +\sum_{n\geq1}\sum_{j=1}^5p_{X,j,n}F_j(n,z)
\end{align}
satisfies $\mathcal P_X'(z)=F_X(z)$ and has zero constant term
at the singular origin.  For real $b>0$,
\[
 \operatorname{FP}\int_0^bF_X(x)\,dx=\mathcal P_X(b).
\]
For $0<a<b$ the ordinary integral is
$\mathcal P_X(b)-\mathcal P_X(a)$.
\end{corollary}
\begin{proof}
Each $F_j$ is the primitive vanishing at zero of
$(n+z)^{-j}-n^{-j}$.  On compact sets it is $O(n^{-j-1})$;
the coefficient bounds from Lemma~\ref{q5:ML} therefore justify normal
convergence and differentiation.  The series converges uniformly on
each disk $|z|\leq r<1$, where all its summands vanish at zero.
The explicitly displayed singular terms and logarithm consequently
leave no constant term.
\end{proof}

\begin{corollary}[All fixed polygamma orders]\label{q5:derivative}
For every integer $r\geq1$, away from the poles,
\begin{equation}\label{q5:derivatives}
 D^rF_X(z)=(-1)^r r!
       \sum_{n\geq0}\sum_{j=1}^5
       \binom{r+j-1}{j-1}\frac{p_{X,j,n}}{(n+z)^{r+j}}.
\end{equation}
The series converges normally.  The left sides are explicitly
\[
 D^rF_5=
 \sum_{r_1+\cdots+r_5=r}\frac{r!}{r_1!\cdots r_5!}
                  \prod_{i=1}^5\psi^{(r_i)},
 \qquad
 D^rF_I=
 \sum_{r_1+r_2+r_3=r}\frac{r!}{r_1!r_2!r_3!}
         \psi^{(r_1)}\psi^{(r_2+1)}\psi^{(r_3+1)}.
\]
At $z=1$ these are finite polynomials in $\gamma$ and ordinary
zeta values, using
$\psi^{(j)}(1)=(-1)^{j+1}j!\zeta(j+1)$ for $j\geq1$.
\end{corollary}
\begin{proof}
Differentiate \eqref{q5:MLformula}.  For $r\geq1$ all constant
subtractions disappear, and the summands are bounded on compact sets
by a constant times $n^{-r-1}(1+\log^4n)$.  The product rule gives
the displayed left sides.
\end{proof}

\subsection{Verification and scope}

The accompanying exact replay checks 87 identities: principal
coefficients, constant Laurent terms, residue differences, both finite
harmonic summation formulas, all finite asymptotic constants,
nonlogarithmic reductions, spectral extraction, endpoint terms,
primitive derivatives, and the derivative kernels through order eight.
The finite harmonic checks are symbolic identities for arbitrary $N$;
they are not a fit to a finite sample of integers.  As regression checks,
the same summation code recovers the previously proved cubic and
quartic constants.

Independent numerical evaluations give
\begin{align*}
 Q_5&=-8.69030396262453826629558966236565357677029354492388429266\ldots,\\
 I&=-1.01946582194172926744369821391554613746742300351148475173\ldots.
\end{align*}
One route integrates the convergent Taylor series of the functions
after their fifth-order singularity has been factored out near zero,
and performs ordinary quadrature away from zero.  Cutoffs $1/8$ and
$1/16$ agree to about 77 decimal digits.  A second route uses
\eqref{q5:Iarithmetic} and \eqref{q5:Qarithmetic}, summing a finite
prefix and expressing a truncated digamma asymptotic tail through
Hurwitz-zeta derivatives.  At the finest recorded settings the
differences from the direct calculation are below
$1.02\cdot10^{-48}$ and $3.83\cdot10^{-48}$, respectively.
These are floating-point cross-checks, not interval-certified
enclosures.  The exact identities rest on the proofs above.

The advance relative to the inspected quartic report is the complete
fifth-power bridge, its mixed-coordinate identification, all finite
corrections, and the normalized primitive and derivative consequences.
The classical partial-fraction and multiple-zeta tools are not claimed
as new.  The scalar-polynomial elimination question remains distinct:
the presently proved cancellation naturally uses the additional
integral $I$ or, equivalently, the differential polynomial
$\mathcal A_5$.
